%% parts/p07/cut-elimination-and-interpolation.tex --- Phase-4 P7.2
%% Migrated source: v1.1 lines 11711--11960. Depth D-P8-003/004 and
%% N-a-069/070. R-012 is homed on sheet:lk-defneg and sheet:lj.

\chapter{Cut Elimination and Interpolation}
\label{sec:p07-cut-elimination-and-interpolation}
\chaptermark{Cut Elimination and Interpolation}

This chapter is scoped to
\hyperref[sheet:lk-defneg]{\(\mathsf{LK}_{\mathrm d}\)} and
\hyperref[sheet:lj]{\(\mathsf{LJ}\)} exactly as registered.  It is not a
theorem about canonical primitive-negation LK, the equality extension,
arithmetic induction, modal or fixed-point rules, or an infinitary calculus.
The Mix rows required by repair R-012 live on precisely these two sheets.
Complete case analyses are available in
\cite{TroelstraSchwichtenberg2000,NegriVonPlato2001}.

\section{Mix measures and elimination}

\begin{hbdefinition}{def:p07-mix-measure}{Formula rank, derivation height, and the Mix measure}
The \emph{formula rank}
\(\operatorname{rk}_{\mathrm{fm}}(\varphi)\) is the number of logical
symbols in \(\varphi\).
The
\emph{derivation height} \(\operatorname{ht}(D)\) is the greatest number of
inference steps on a leaf-to-root branch.  On the two active sheets,
\(\varphi^k\) in a Mix row means exactly \(k\) comma-separated occurrences
of the mix formula \(\varphi\).  If \(D_1,D_2\) are the immediate
subderivations of that
Mix instance, its measure is
\[
 \bigl\langle \operatorname{rk}_{\mathrm{fm}}(\varphi),
 \operatorname{ht}(D_1)+\operatorname{ht}(D_2)\bigr\rangle,
\]
ordered lexicographically.  Only after a Cut has been read as the special Mix
instance with multiplicities one do we call these coordinates its cut rank
and cut height.  Rank measures formula syntax; height measures a proof tree;
the distinction is required by the two independent decreases below.
\index{formula rank}
\index{derivation height}
\index{Mix measure}
\end{hbdefinition}

\begin{hblemma}{thm:p07-cut-mix-equivalence}{proved=annex:p7.cut.1}{Equipollence of Cut and Mix}
Replacing Cut by the registered Mix rule does not change the derivable
endsequents of \(\mathsf{LK}_{\mathrm d}\) or \(\mathsf{LJ}\).
Consequently, a Mix-free derivation in the Mix formulation is cut-free, and
an elimination procedure for Mix eliminates Cut.
\end{hblemma}

\begin{hbmethod}{def:p07-cut-elimination-method}{Cut-elimination method}
\begin{description}
  \item[Input.] A finite \(\mathsf{LK}_{\mathrm d}\)- or
        \(\mathsf{LJ}\)-derivation.
  \item[Normalize the rule surface.] Read every Cut as the registered Mix
        instance with multiplicities one.
  \item[Reduce.] Choose an uppermost Mix of maximal formula rank.  Commute a
        nonprincipal Mix above the final inference; delete weakening and
        initial cases; absorb contraction by increasing the appropriate Mix
        multiplicity; replace a principal/principal instance by Mixes on
        proper subformulas plus equal-rank residual Mixes of smaller height;
        and alpha-rename eigenvariables before quantifier permutations and
        capture-avoiding substitutions.
  \item[Measure.] Iterate in the well-founded lexicographic order of
        Definition~\ref{def:p07-mix-measure}.
  \item[Output.] A derivation in the same calculus, of the same endsequent,
        containing no Mix and hence no Cut.
\end{description}
The name abbreviates this calculus-specific transformation, not an appeal to
semantic completeness.  Every listed family is used by the theorem; the
omitted per-rule table is recorded in the annex.
\end{hbmethod}

\begin{hbtheorem}{thm:p07-cut-elimination}{sketched=annex:p7.cut.2}{Cut elimination for defined-negation LK and LJ}
Every \(\mathsf{LK}_{\mathrm d}\)-derivation can be transformed into a
cut-free \(\mathsf{LK}_{\mathrm d}\)-derivation of the same endsequent.
Every \(\mathsf{LJ}\)-derivation can be transformed into a cut-free
\(\mathsf{LJ}\)-derivation of the same endsequent.
\end{hbtheorem}

\begin{proofboundary}{def:p07-cut-proof-boundary}{What the inherited sketch still owes}
The annex records the six load-bearing reduction families and the decreasing
measure.  A complete proof would additionally give every commuting and
principal reduction for every logical and structural rule on both sheets,
including contraction multiplicities and all quantifier/eigenvariable cases.
Those omissions are why the theorem remains \emph{sketched}.  It cannot be
transferred to \texttt{sheet:lk}, \texttt{sheet:lk-eq}, or another calculus
merely by calling a derivation cut-free.
\end{proofboundary}

\section{Analytic derivations and proof search}

\begin{hbdefinition}{def:p07-analytic-derivation}{Analytic derivation in the active calculi}
A propositional \(\mathsf{LK}_{\mathrm d}\)- or \(\mathsf{LJ}\)-derivation
is \emph{analytic} iff it is cut-free and every formula occurrence in it is
a literal subformula of a formula in its endsequent.  A first-order
derivation is \emph{analytic modulo terms} iff it is cut-free and every
formula occurrence is a subformula of an endsequent formula after term
substitution and eigenvariable alpha-renaming.  The qualification is
nonredundant: quantifier rules introduce instances
\(\varphi[t/x]\) that need not
be literal subformulas.  These definitions are used to separate the
propositional theorem from the first-order warning.
\index{analytic derivation}
\index{subformula property}
\end{hbdefinition}

\begin{hblemma}{thm:p07-cut-free-subformula}{proved=annex:p7.cut.3}{Subformula property of cut-free propositional derivations}
Every cut-free propositional \(\mathsf{LK}_{\mathrm d}\)- or
\(\mathsf{LJ}\)-derivation has the literal subformula property and hence is
analytic in the sense of Definition~\ref{def:p07-analytic-derivation}.
\end{hblemma}

\begin{hbcorollary}{thm:p07-propositional-analytic-existence}{sketched=annex:p7.cut.4}{Propositional analytic existence}
\par\noindent
Every derivable propositional \(\mathsf{LK}_{\mathrm d}\)- or
\(\mathsf{LJ}\)-sequent has an analytic derivation of the same endsequent.
\end{hbcorollary}

\begin{hbcorollary}{thm:p07-empty-sequent}{sketched=annex:p7.cut.5}{Underivability of the empty sequent}
The empty sequent \(\varnothing\Rightarrow\varnothing\) has no
\(\mathsf{LK}_{\mathrm d}\)-derivation.
\end{hbcorollary}

\begin{hbwarning}{def:p07-cut-free-search}{Cut-free does not mean efficient or finite search}
Cut elimination can increase proof size nonelementarily.  Even after Cut is
removed, first-order backward search can branch over unbounded term choices.
The literal subformula property therefore supports a propositional decision
argument only after the finite rule-search space of the exact calculus is
verified; it does not directly decide first-order \(\mathsf{LJ}\).
\end{hbwarning}

\section{Midsequents and the Herbrand route}

\begin{hbdefinition}{def:p07-regular-cut-free-fol-derivation}{Regular cut-free derivation}
\par\noindent
For this section, a \emph{regular derivation} is a
finite \(\mathsf{LK}_{\mathrm d}\)-derivation that contains no Cut, Mix,
equality, or induction inference and assigns distinct eigenvariables to
distinct quantifier inferences, each fresh for the conclusion required by its
rule.  This is a metalinguistic classification of derivation trees, not a new
calculus.  Regularity is needed for the quantifier permutations in the next
theorem and differs from cut-freeness by its global eigenvariable condition.
\index{regular derivation}
\index{midsequent theorem}
\end{hbdefinition}

\begin{hbtheorem}{thm:p07-midsequent}{cited=Takeuti1987}{Midsequent theorem, equality-free prenex form}
\par\noindent
If a prenex endsequent has a regular equality-free cut-free derivation on
\(\mathsf{LK}_{\mathrm d}\), that derivation can be permuted so that no
propositional logical inference occurs below a quantifier inference, with
tree roots drawn at the bottom.  A quantifier-free sequent then separates the
upper propositional part from the lower quantifier part.
\end{hbtheorem}

\begin{hbtheorem}{thm:p07-herbrand-midsequent}{cited=Takeuti1987}{Herbrand route from a midsequent}
\par\noindent
Let
\[
 \exists x_1\cdots\exists x_k\,\varphi(x_1,\ldots,x_k)
\]
be a valid equality-free prenex sentence with quantifier-free matrix, in a
signature containing a ground term or after adding one fresh constant.  A cut-free
derivation followed by the midsequent permutation exposes finitely many
ground tuples \((t_1^j,\ldots,t_k^j)\), \(1\leq j\leq n\), for which
\[
 \varphi(t_1^1,\ldots,t_k^1)\lor\cdots\lor
 \varphi(t_1^n,\ldots,t_k^n)
\]
is propositionally valid.
\end{hbtheorem}

This is a proof-theoretic route to the citation-only validity form in
Theorem~\ref{thm:p02-herbrand-validity}; it neither upgrades that theorem nor
replaces its finite-clause unsatisfiability form or the epsilon-calculus
route.  An equality-bearing Herbrand theorem would require equality instances
or an equality calculus and is not asserted.

\section{Maehara interpolation and definability}

\begin{hbdefinition}{def:p07-partition-interpolant}{Partition and partition interpolant}
Let an equality-free endsequent of sentences be written
\(\Gamma_1,\Gamma_2\Rightarrow\Delta_1,\Delta_2\).  The ordered four-tuple of
displayed subcontexts is a \emph{partition} of the sequent.  A
\emph{partition interpolant} is a sentence \(\chi\) such that
\(\Gamma_1\Rightarrow\Delta_1,\chi\) and
\(\chi,\Gamma_2\Rightarrow\Delta_2\) are derivable in
\(\mathsf{LK}_{\mathrm d}\), and every nonlogical
predicate, function, and constant symbol in \(\chi\) occurs both in a formula
of \(\Gamma_1\cup\Delta_1\) and in a formula of
\(\Gamma_2\cup\Delta_2\).  Here \(\bot\) is logical and equality is absent.

The term is necessary because an arbitrary intermediate consequence lacks
the common-vocabulary clause; it names a sentence relative to a partition,
not a new consequence relation.  That clause is used to obtain Craig's
vocabulary restriction below.
\index{partition interpolant}
\end{hbdefinition}

\begin{hbmethod}{def:p07-maehara-method}{Maehara's partition method}
Induct on a supplied cut-free equality-free
\(\mathsf{LK}_{\mathrm d}\)-derivation together with a partition of its
endsequent.  Construct an interpolant at initial and \(\bot L\) leaves;
propagate it through weakening and contraction; combine premise interpolants
separately for every placement of each propositional principal formula; and
in the four quantifier cases quantify away eigenvariables not in the common
language while preserving the two required derivations.  Applied to a
cut-free derivation of \(\varphi\Rightarrow\psi\) with partition
\((\varphi;\varnothing\mid\varnothing;\psi)\), the output has the common
nonlogical vocabulary needed for Craig interpolation.  The annex records why
the omitted rule-by-partition table leaves the local lemma a sketch.
\end{hbmethod}

\begin{hblemma}{thm:p07-maehara-partition}{sketched=annex:p7.interp.1}{Maehara's partition lemma, sentence form}
Every partition of the endsequent of a cut-free equality-free
\(\mathsf{LK}_{\mathrm d}\)-derivation, when all four endsequent contexts
contain sentences, has a partition interpolant in the sense of
Definition~\ref{def:p07-partition-interpolant}.
\end{hblemma}

\begin{hbtheorem}{thm:p07-craig-interpolation}{cited=Craig1957}{Craig interpolation}
Let \(\varphi,\psi\) be first-order sentences and suppose
\(\entails\varphi\to\psi\).  Then there is a sentence \(\chi\) such that
\(\entails\varphi\to\chi\), \(\entails\chi\to\psi\), and every
nonlogical symbol of \(\chi\) occurs in both \(\varphi\) and \(\psi\).
Equality and \(\bot\) may occur in \(\chi\) regardless.  Degenerate cases
are included because \(\bot\) and \(\neg\bot\) are available when
\(\varphi\) is unsatisfiable or \(\psi\) is valid.
\end{hbtheorem}

\begin{hbtheorem}{thm:p07-beth-definability}{cited=Beth1953}{Beth definability}
Let \(T\) be a first-order theory in a signature
\(\Lang\cup\{R\}\), where \(R\) is \(n\)-ary.  Suppose any two models of
\(T\) with the same \(\Lang\)-reduct interpret \(R\) identically.  Then
there is a \(\Lang\)-formula
\(\chi(x_1,\ldots,x_n)\) such that
\[
 T\derives_{\mathsf{ND}_{\mathsf{FOL}}}
 \forall x_1\cdots\forall x_n\,
 \bigl(R(x_1,\ldots,x_n)\leftrightarrow
       \chi(x_1,\ldots,x_n)\bigr).
\]
\end{hbtheorem}

The standard Beth route renames \(R\), compares two copies of the theory,
applies compactness and interpolation, and removes the renamed symbol;
\cite{Hodges1993} supplies that argument.  Robinson joint consistency is
stated separately in Part~2.  Neither cited source theorem receives an annex
item: the new local Maehara sketch supplies a method entry, not an upgrade of
the inherited evidence classes for Craig or Beth.

\begin{proofboundary}{def:p07-interpolation-proof-boundary}{Conditional local method versus cited theorems}
\par\noindent
The local Maehara entry assumes an already supplied cut-free derivation.
Obtaining one from an arbitrary derivation loads the sketched cut-elimination
theorem; obtaining a derivation from semantic consequence loads the exact
completeness theorem and its hypotheses.  A complete Maehara proof would also
need every initial, structural, propositional-placement, and quantifier case.
The register-driven method therefore remains sketched, while the migrated
Craig and Beth results remain cited exactly as in v1.1.
\end{proofboundary}

\hbstub{D-P8-004 remains open at its E ceiling: completing it requires a
fully checked Maehara induction through every rule and partition case, followed
by a locally supported Craig proof; the present partition lemma is explicitly
sketched and the inherited Craig theorem remains cited.}

\hbstub{N-a-070 remains open at its E ceiling for the same evidence reason:
the handbook supplies the Maehara method and exact cited Craig and Beth
statements, but no complete local Craig proof or complete Beth-corollary
derivation whose status could honestly be upgraded.}

\paragraph{Source disposition.}
The source exercise requesting the principal reduction for a conjunction Mix
is retired as an exercise; its smaller-rank obligation is retained in
Method~\ref{def:p07-cut-elimination-method} and in the annex omission ledger.
The exercise comparing propositional and first-order proof search is retired
with its distinction preserved in Warning~\ref{def:p07-cut-free-search}.
No solution to either prompt is reclassified as a proof.
